% ECONOMICS_PROBLEM: {"id": "G28-112", "title": "Nonbank systemic regulation", "jel_code": "G28", "source_id": "OP-0325"}
% !TeX program = xelatex
\documentclass[11pt,a4paper]{article}
\usepackage[margin=23mm]{geometry}
\usepackage{fontspec}
\setmainfont{Latin Modern Roman}
\usepackage{amsmath,amssymb,mathtools}
\usepackage{xcolor,enumitem,booktabs,longtable,array,xurl,hyperref}
\definecolor{muted}{HTML}{53636C}
\hypersetup{colorlinks=true,urlcolor=blue,linkcolor=blue}
\setlength{\parindent}{0pt}
\setlength{\parskip}{5pt}
\newcommand{\R}{\mathbb R}
\newcommand{\E}{\mathbb E}
\newcommand{\N}{\mathbb N}
\newcommand{\ind}{\mathbf 1}
\newcommand{\Var}{\operatorname{Var}}
\newcommand{\Cov}{\operatorname{Cov}}
\newcommand{\supp}{\operatorname{supp}}
\DeclareMathOperator*{\argmin}{arg\,min}
\DeclareMathOperator*{\argmax}{arg\,max}
\newcommand{\entrymeta}[3]{{\sffamily\small\color{muted}\textbf{\texttt{#1}}\quad|\quad #2\par #3\par}}
\newcommand{\Problemref}[1]{\texttt{#1}}
\newcommand{\JELref}[2]{\texttt{#2} (JEL \texttt{#1})}
\begin{document}
\section*{Nonbank systemic regulation}
\textbf{Problem ID:} \texttt{G28-112}\quad \textbf{JEL:} \texttt{G28}\par
% BEGIN ECONOMICS STATEMENT
\label{problem:OP-0325}
{\sffamily\small\color{muted}Original subject group: Banking, credit and financial stability\par}
\label{definitions:problem:30007607}
\textbf{Audit classification:} Explicit published open question. The BIS primary open-question section explicitly addresses the balance of prudential costs and stability benefits. Exact tool sets, loss functions and robust probabilities are editorial.
\subsection*{Definitions and precise research target}
Consider nonbank intermediaries funding illiquid assets with redeemable claims, secured borrowing or derivatives. A feasible regulation \(a\) specifies leverage limits, redemption terms, swing pricing and initial or variation-margin rules. Investors choose leverage and liquidity buffers before shocks and trade or redeem after them; prices clear with explicitly modeled market impact. Let \(L(a,\omega)\) measure default, fire-sale and liquidity externalities in shock state \(\omega\), and \(C(a)\) routine intermediation and compliance costs. For a specified uncertainty set \(\mathcal P\), the design task is
\[J(a)=C(a)+\sup_{P\in\mathcal P}E_P[L(a,\omega)],\qquad J_*=\inf_{a\in A}J(a).\]
The set \(A\) includes legal feasibility and operational implementability restrictions. Characterize how optimal tools depend on redemption mismatch, cross-institution exposures and endogenous margin spirals. Identify the information needed to estimate losses and costs, and test whether activity shifts to uncovered institutions when constraints tighten. Include changes in ordinary financing capacity, rather than treating regulation as costless. Report whether tools reduce individual-institution losses while increasing system-wide externalities. Existence of an optimal rule for a fixed shock law is a model result; robust implementation and empirical calibration remain separate parts of the research task.

\textbf{Additional formal definitions and scope (editorial).} Specify each institution's asset maturity and redemption liabilities, leverage measurement including derivative exposures, and the margin update equations. Swing pricing adjusts a redeeming investor's payment by a stated function of transaction costs; it is not identical to a redemption gate. Variation margin settles current mark-to-market losses, whereas initial margin is posted against prospective exposure. The regulator observes a declared information set, and policies are measurable with respect to it. The optimization problem includes endogenous leverage and redemption choices and a selected market-clearing allocation. Nonempty compact \(A\) and a finite lower semicontinuous loss objective give existence of a minimum by standard compactness; finding justified primitives, tool interactions and their identified welfare consequences is the research agenda. For a nonempty feasible set and a finite optimal value, the design target is an \(\varepsilon\)-optimal feasible rule for each specified \(\varepsilon>0\): its cost is at most the infimum plus \(\varepsilon\), or its welfare is at least the supremum minus \(\varepsilon\). Report infeasibility or an unbounded value separately. An exact optimizer is requested only under additional attainment assumptions; it need not exist for the general rule class.
\subsection*{Variants and assumption changes}
\begin{enumerate}
\item \textbf{Single institution class (editorial).} Restrict to redeemable funds with observed market impact and compare swing pricing, gates and leverage rules.
\item \textbf{Activity-wide regulation (source-motivated extension).} Include funds, secured lenders and derivatives together, allowing margin feedback and migration beyond the regulatory perimeter.
\end{enumerate}
\subsection*{References and source audit}
\begin{enumerate}
\item BIS WP 972 cover verifies authors; October 2021 original, January 27, 2022 revision inspected; no DOI assigned. Locator: Section 5, printed pp.38--39 (PDF pp.41--42). Support: Nonbank regulation/backstop agenda. Full January 2022 revision accessible. URL: \url{https://www.bis.org/publications/working-paper-972-non-bank-financial-intermediaries-and-financial-stability.pdf}.
\item Official January 9, 2026 web version; locator: Prudential Architecture / Institutional investors, regulatory-perimeter paragraph. Broad agenda; exact model editorial. URL: \url{https://www.bankofengland.co.uk/research/bank-of-england-agenda-for-research}.
\end{enumerate}
\begin{enumerate}\raggedright
\item\hypertarget{definitions:source:30007607:1}{} Sirio Aramonte; Andreas Schrimpf; Hyun Song Shin. 2021. \emph{Non-bank financial intermediaries and financial stability}. Section 5, printed pp.38--39 (PDF pp.41--42).
\url{https://www.bis.org/publications/working-paper-972-non-bank-financial-intermediaries-and-financial-stability.pdf}.
\item\hypertarget{definitions:source:30007607:2}{} Bank of England. 2026. \emph{Bank of England Agenda for Research, 2025–2028}. Prudential Architecture / Institutional investors, regulatory-perimeter paragraph.
\url{https://www.bankofengland.co.uk/research/bank-of-england-agenda-for-research}.
\end{enumerate}

\subsection*{Common conventions: Source mathematical conventions}

These conventions apply to each entry unless explicitly replaced.

\textbf{Models and identification.} A primitive model \(m\) specifies all probability laws, feasible choices, information, equilibrium and measurement rules used by its target. The admissible class \(\mathcal M\) is fixed independently of outcomes. If \(P_O(m)\) is its observable probability law and \(\tau(m)\) the target, its identified set at observable law \(P\) is
\[
 \mathcal I_\tau(P)=\{\tau(m):m\in\mathcal M,\ P_O(m)=P\}.
\]
Point identification means that this set is a singleton. An empty set signals incompatibility of data and assumptions. Coordinate bounds need not describe the joint set. Bounds are sharp only if every retained target is attainable or arbitrarily approachable by compatible models. Finite bounds require explicit restrictions on unobserved quantities; unrestricted confounding, hidden wealth, extrapolation or arbitrary equilibrium selection can yield uninformative sets.

\textbf{Causal interventions.} A potential outcome \(Y_i(p)\) is the outcome under a complete policy regime \(p\), including timing, financing, information and market spillovers. The target population and units are fixed before treatment. With interference, use the complete assignment vector or a justified exposure mapping; individual no-interference notation alone cannot identify an economy-wide intervention. A difference of mean potential outcomes does not require the cross-world joint distribution. Conditioning on post-treatment migration, survival, enrollment or employment changes the estimand and needs separate justification.

\textbf{Statistical statements.} An estimator, confidence set, forecast or test specifies a sampling law, model class and loss/coverage criterion. Uniform \((1-\alpha)\)-coverage means \(\inf_{m\in\mathcal M}\Pr_m\{\tau(m)\in C_n\}\geq1-\alpha\), or an explicitly asymptotic version. The whole parameter space trivially has coverage, so informativeness requires a stated width, risk or power objective. A forecasting improvement is relative to a named benchmark, information set and evaluation population.

\textbf{Equilibrium and optimization.} An equilibrium comprises feasible plans, optimality, beliefs where needed, budget constraints and all market/resource clearing. The primitives or a stated benchmark define each correspondence \(\mathcal E(p;m)\). Nonemptiness is assumed or proved, never inferred just from compact policy sets. A selected-equilibrium target uses a declared selection rule; a robust target quantifies over all permitted equilibria. Use suprema or infima unless attainment is established. An \(\varepsilon\)-optimal policy is feasible and within \(\varepsilon\) of the finite optimum value. Report an empty feasible set separately.

\textbf{Welfare and accounting.} Normative utility, interpersonal normalization, social weights, numeraire, discounting and the use of public receipts are primitives. Behavioral utility need not coincide with the welfare criterion. Household welfare includes ownership income and rents once; adding profits again to compensating variation generally double-counts them. A monetary equivalent is defined by an explicit transfer and price convention, with monotonicity and range conditions. Average effects, mechanism contributions, transfers and welfare losses are different targets. Infinite sums require convergence and dynamic feasible plans require terminal or no-Ponzi conditions.

\textbf{Source versions.} A working-paper date, revision date and journal publication year are distinguished. A DOI for a journal version is not automatically the DOI of an earlier manuscript. Locators refer to the stated version and printed pages where specified. A metadata-only check does not verify a claim about a full-text conclusion. Every entry's source subsection narrows the provenance claim accordingly.
% END ECONOMICS STATEMENT
\end{document}
